%%%%%%%%%%%%%%%%%%%%%%%%%%%%%%%%%%%%%%%%%%%%%%%%%%%%%%%%%%%%%%%%%%%%%%
% How to use writeLaTeX: 
%
% You edit the source code here on the left, and the preview on the
% right shows you the result within a few seconds.
%
% Bookmark this page and share the URL with your co-authors. They can
% edit at the same time!
%
% You can upload figures, bibliographies, custom classes and
% styles using the files menu.
%
%%%%%%%%%%%%%%%%%%%%%%%%%%%%%%%%%%%%%%%%%%%%%%%%%%%%%%%%%%%%%%%%%%%%%%

\documentclass[12pt]{article}

% `inputenc` não deve ser usado com fontspec (XeTeX/LuaTeX); com pdfLaTeX mantém UTF-8.
\usepackage{iftex}
\ifPDFTeX
  \usepackage[utf8]{inputenc}
\fi

\usepackage{sbc-template}
\usepackage{graphicx,url}
\usepackage{float} 
\usepackage[brazil]{babel}
\usepackage{tabularx}
\newcolumntype{C}{>{\centering\arraybackslash}X} % <-- modified
\usepackage{amsmath}
\usepackage{pgfplots}
\pgfplotsset{compat=1.18, every axis/.append style={font=\footnotesize}}
\usetikzlibrary{babel}
\usepackage[inline]{enumitem}
\usepackage{tcolorbox}


%%%%%%%%%%%%%%%%% Pacotes para inserção de algoritmos - BEGIN
\usepackage[ruled,lined]{algorithm2e}


% Pacote para a definição de novas cores
\usepackage{xcolor}
% Definindo novas cores
\definecolor{verde}{rgb}{0.25,0.5,0.35}
\definecolor{jpurple}{rgb}{0.5,0,0.35}
\definecolor{darkgreen}{rgb}{0.0, 0.2, 0.13}
%\definecolor{oldmauve}{rgb}{0.4, 0.19, 0.28}

\usepackage{listings}
\lstset{
    language=C++,
    basicstyle=\ttfamily\small,
    keywordstyle=\color{jpurple}\bfseries,
    stringstyle=\color{red},
    commentstyle=\color{verde},
    morecomment=[s][\color{blue}]{/**}{*/},
    extendedchars=true,
    showspaces=false,
    showstringspaces=false,
    numbers=left,
    numberstyle=\tiny,
    breaklines=true,
    backgroundcolor=\color{white!10},
    breakautoindent=true,
    captionpos=b,
    xleftmargin=0pt,
    tabsize=2
}
%%%%%%%%%%%%%%%%% Pacotes para inserção de algoritmos - END

     
\sloppy

\title{Análise Empírica de Estruturas Union--Find para Conectividade Dinâmica}

\author{Lucas Picinin Campos Lutti\inst{1} }


\address{%
Pontifícia Universidade Católica de Minas Gerais (PUC Minas)\\
Curso de Engenharia de Software
  \email{lucaspicinin@hotmail.com}
}

\begin{document} 

\maketitle

\begin{abstract}
  This paper compares three Disjoint Set Union (Union--Find) variants: a naive implementation,
  union by rank, and union by rank combined with path compression. The comparison is based on
  reproducible sequences of \texttt{Find} and \texttt{Union} operations, with measurements of total
  execution time and instrumented accesses to the internal \texttt{parent} and \texttt{rank} arrays.
  The results highlight the practical performance gap between the naive structure and the
  classical optimizations associated with Tarjan-style analysis.
\end{abstract}
     
\begin{resumo} 
  Este trabalho compara três variantes de conjuntos disjuntos (\emph{Union--Find}): uma implementação
  ingênua, outra com união por rank e uma terceira com união por rank associada à compressão de caminho.
  A comparação é feita sobre sequências reproduzíveis de operações \texttt{Find} e \texttt{Union}, medindo
  o tempo total de execução e os acessos instrumentados aos vetores internos da estrutura.
  Os resultados mostram, na prática, a diferença de desempenho entre a variante ingênua e as
  heurísticas clássicas analisadas por Tarjan.
\end{resumo}


\begin{tcolorbox}
\footnotesize
\textbf{Fundamentos de Projeto e Análise de Algoritmos (FPAA) --- Trabalho Prático}\\
PUC Minas, Engenharia de Software. Os dados do ambiente de experimentos estão descritos na Metodologia e em \texttt{docs/AMBIENTE.md} no repositório do projeto.
\end{tcolorbox}


\section{Introdução}
\label{sec:introducao}

Problemas de conectividade aparecem em diferentes contextos da computação sempre que é necessário manter grupos de elementos e responder, com rapidez, se dois deles pertencem ao mesmo conjunto. No caso de grafos, essa situação corresponde à conectividade dinâmica incremental: novas arestas são inseridas e deseja-se saber se dois vértices passaram a fazer parte da mesma componente conexa.

A estrutura de dados \emph{Disjoint Set Union} (DSU), também chamada de \emph{Union--Find}, oferece as operações \texttt{Find} (obtenção do representante de um conjunto) e \texttt{Union} (fusão entre conjuntos distintos). Quando implementada de forma ingênua, essa estrutura tende a perder desempenho à medida que o número de operações cresce. Já heurísticas clássicas, como união por rank e compressão de caminho, reduzem substancialmente esse custo e levam à análise amortizada associada a Tarjan, expressa em termos de $\mathcal{O}(\alpha(n))$, em que $\alpha$ é a inversa da função de Ackermann \cite{tarjan1975efficiency}.

Neste trabalho, comparam-se três variantes de DSU: uma implementação ingênua, uma com união por rank e outra com união por rank combinada à compressão de caminho. Para essa comparação, o código foi organizado de forma modular, instrumentado para registrar leituras e escritas nos vetores \texttt{parent} e \texttt{rank}, e executado sobre sequências reproduzíveis de operações. A análise de Tarjan \cite{tarjan1975efficiency} sustenta a interpretação teórica das heurísticas estudadas, enquanto os tempos, acessos e gráficos apresentados ao longo do texto vêm do benchmark descrito na Seção~\ref{sec:metodologia}. Como referência de aplicação, menciona-se o algoritmo de Kruskal, que utiliza a mesma interface de \emph{Union--Find} para detectar ciclos durante a construção de árvores geradoras mínimas.

As seções seguintes apresentam a fundamentação teórica, o posicionamento do trabalho, a metodologia adotada, a discussão dos resultados e, por fim, as conclusões do estudo.

\section{Fundamentação Teórica}
\label{sec:fundamentacao}

\subsection{Conjuntos disjuntos e operações}

Seja um universo $\{0,1,\ldots,n-1\}$ particionado em conjuntos disjuntos. Cada conjunto é representado por um elemento identificador (raiz). A operação $\texttt{Find}(x)$ devolve o representante do conjunto que contém $x$; $\texttt{Union}(a,b)$ funde os conjuntos que contêm $a$ e $b$, caso sejam distintos. Duas consultas $\texttt{Find}(x)$ e $\texttt{Find}(y)$ retornam o mesmo valor se, e somente se, $x$ e $y$ pertencem ao mesmo conjunto.

\subsection{Floresta com união ingênua}

Uma forma clássica de representar conjuntos disjuntos consiste em manter uma floresta por meio de um vetor $\texttt{parent}$, no qual cada elemento aponta para seu pai e toda raiz satisfaz $\texttt{parent}[r]=r$. A operação $\texttt{Find}$ percorre esses ponteiros até alcançar a raiz, enquanto $\texttt{Union}$ localiza as raízes de $a$ e $b$ e faz uma delas apontar para a outra. Sem qualquer critério de balanceamento, a altura da árvore pode crescer linearmente, levando $\texttt{Find}$ a custo $\mathcal{O}(n)$ no pior caso e, portanto, sequências de $m$ operações podem atingir $\mathcal{O}(mn)$.

\subsection{União por rank}

Na heurística de união por rank, associa-se a cada nó um \emph{rank}, que funciona como majorante da altura da subárvore enraizada naquele ponto. Em $\texttt{Union}$, a raiz de menor rank é anexada à de maior rank; em caso de empate, incrementa-se o rank da nova raiz. Com essa estratégia, a altura das árvores permanece em $\mathcal{O}(\log n)$, de modo que cada $\texttt{Find}$ ou $\texttt{Union}$ custa $\mathcal{O}(\log n)$ no pior caso.

\subsection{Compressão de caminho e resultado de Tarjan}

Na compressão de caminho, cada nó visitado durante $\texttt{Find}$ passa a apontar diretamente para a raiz do conjunto, o que achata a estrutura ao longo do tempo. Quando essa heurística é combinada à união por rank, a análise amortizada sobre uma sequência de $m$ operações (com $m \geq n$) mostra que o tempo total é $\mathcal{O}(m\,\alpha(n))$, onde $\alpha(n)$ é a inversa de Ackermann, função que cresce tão lentamente que, para todos os valores de $n$ encontrados em aplicações práticas, $\alpha(n)\leq 4$ \cite{tarjan1975efficiency}.

\subsection{Interpretação prática de $\alpha(n)$}

A função de Ackermann cresce muito rapidamente, enquanto sua inversa $\alpha(n)$ cresce de forma extremamente lenta. Por isso, embora a complexidade $\mathcal{O}(\alpha(n))$ não seja constante em sentido estrito, na prática as operações se comportam quase como tempo constante amortizado para tamanhos de instância usuais. Esse comportamento justifica a presença recorrente do DSU em algoritmos de grafos, como Kruskal.

\section{Trabalhos Relacionados}
\label{sec:relacionados}

Estruturas do tipo DSU aparecem com frequência em problemas de conectividade incremental e em algoritmos clássicos como o de Kruskal para árvores geradoras mínimas. Nesse contexto, o trabalho de Tarjan estabelece a análise amortizada da combinação entre união por rank e compressão de caminho, obtendo o limite $\mathcal{O}(m\,\alpha(n))$ para uma sequência de $m$ operações \cite{tarjan1975efficiency}.

Em materiais didáticos e implementações voltadas a estudo, é comum comparar versões sem heurísticas com variantes mais otimizadas para evidenciar a diferença entre seus custos. O foco deste trabalho está nessa comparação, mas com um protocolo experimental explícito: as três variantes são avaliadas sobre as mesmas sequências reproduzíveis, com medição de tempo e de acessos à memória, para aproximar a análise assintótica do comportamento observado na execução.

\section{Metodologia}
\label{sec:metodologia}

\subsection{Organização do código}

O código foi dividido em três módulos principais de DSU em C++17, sem bibliotecas prontas para a lógica de conjuntos disjuntos: \texttt{DsuNaive}, \texttt{DsuUnionRank} e \texttt{DsuTarjan}, organizados em arquivos separados nos diretórios \texttt{include/} e \texttt{src/}. Além deles, o projeto mantém componentes compartilhados com funções bem definidas: \texttt{operacoes} gera as sequências reproduzíveis de teste, \texttt{estatisticas\_acesso} concentra a instrumentação de leituras e escritas, e \texttt{experimentos} executa o mesmo protocolo de medição para todas as variantes. Essa separação foi adotada para que uma mudança em uma implementação de DSU não exigisse alterar o gerador das operações nem a rotina de coleta das métricas, preservando a comparabilidade entre os experimentos.

\subsection{Instrumentação}

As leituras e escritas nos vetores internos da estrutura são contabilizadas por funções auxiliares que atualizam contadores por instância. Na variante ingênua, apenas o vetor \texttt{parent} entra nessa contagem; nas demais, contam-se acessos a \texttt{parent} e a \texttt{rank}. As inicializações realizadas antes da medição ficam fora desses totais, de modo que os valores reportados reflitam apenas o custo da sequência de operações \texttt{Find} e \texttt{Union}. O tempo de execução, por sua vez, é medido com \texttt{std::chrono::steady\_clock} ao redor do laço que reaplica a mesma sequência para cada variante.

\subsection{Cenários experimentais}

Para cada valor de $n$ na lista $\{500, 1000, 2000, 4000, 8000, 16000, 32000, 64000\}$, é gerada uma sequência com $m=20n$ operações por meio de \texttt{std::mt19937}. A fração de operações \texttt{Union} foi fixada em $0{,}25$, enquanto as demais correspondem a chamadas de \texttt{Find} em um único índice. Os índices são sorteados uniformemente em $\{0,\ldots,n-1\}$, e a semente usada para o $i$-ésimo valor de $n$ é $42 + 7919i$ (com $i=0$ para $n=500$), exatamente como no modo \texttt{--all} de \texttt{src/experimentos.cpp}. A mesma sequência é então reaplicada, na mesma ordem, às três variantes. Antes das medições, executa-se um aquecimento (\emph{warm-up}) com \texttt{DsuTarjan} sobre essa sequência, sem registrar seu tempo no CSV.

\subsection{Ambiente e compilação}

Os experimentos foram executados em um computador com processador Apple M1, 8\,GB de memória RAM e sistema operacional macOS 26.3 (build 25D125). A compilação foi feita com \texttt{g++}, correspondente ao Apple clang 17.0.0 (clang-1700.0.13.5), usando as flags \texttt{-std=c++17 -O3 -DNDEBUG}. O executável \texttt{bin/benchmark} grava os dados brutos em \texttt{results/benchmark.csv}. Em seguida, o script \texttt{scripts/export\_latex\_plot\_data.py} converte esses dados para \texttt{artigo/data/tempo\_por\_n.dat}, \texttt{artigo/data/acessos\_por\_n.dat} e \texttt{artigo/data/tabela\_resumo.inc.tex}, que são lidos diretamente pelo \LaTeX{} com \texttt{pgfplots}. As mesmas informações de ambiente ficam registradas em \texttt{docs/AMBIENTE.md} no repositório do projeto.

\section{Discussão dos Resultados}
\label{sec:resultados}

\subsection{Tempo de execução}

O gráfico da Figura~\ref{fig:tempo} foi construído diretamente a partir dos dados exportados de \texttt{results/benchmark.csv} para \texttt{data/tempo\_por\_n.dat} por meio de \texttt{scripts/export\_latex\_plot\_data.py}. Nele, a curva da variante ingênua cresce bem mais rapidamente à medida que $n$ aumenta, o que é compatível com encadeamentos longos em \texttt{Find}. A variante com união por rank já reduz esse crescimento de forma visível, enquanto a combinação com compressão de caminho permanece como a mais rápida nas maiores instâncias testadas, como se espera de um custo amortizado muito baixo por operação \cite{tarjan1975efficiency}.

\begin{figure}[ht]
\centering
\begin{tikzpicture}
\begin{axis}[
  width=0.88\textwidth,
  height=6.8cm,
  xlabel={$n$},
  ylabel={Tempo total (ms)},
  xmode=log,
  log basis x=2,
  ymode=log,
  legend style={at={(0.03,0.97)},anchor=north west,font=\footnotesize},
  grid=major,
]
\addplot+[thick,mark=*,mark size=1.1pt,color=red!75!black] table[x=n,y=naive_ms,col sep=tab] {data/tempo_por_n.dat};
\addplot+[thick,mark=square*,mark size=1.1pt,color=blue!70!black] table[x=n,y=union_rank_ms,col sep=tab] {data/tempo_por_n.dat};
\addplot+[thick,mark=triangle*,mark size=1.3pt,color=green!50!black] table[x=n,y=tarjan_ms,col sep=tab] {data/tempo_por_n.dat};
\legend{Ingênua,União por rank,Rank + compressão}
\end{axis}
\end{tikzpicture}
\caption{Tempo total para processar $m=20n$ operações (25\% \texttt{Union}, 75\% \texttt{Find}) em função de $n$. Fonte: medições do \texttt{benchmark} descrito na Seção~\ref{sec:metodologia}.}
\label{fig:tempo}
\end{figure}

\subsection{Acessos à memória}

Na Figura~\ref{fig:acessos}, vê-se o total de leituras e escritas instrumentadas para a mesma sequência de operações: somente no vetor \texttt{parent} na variante ingênua e em \texttt{parent} mais \texttt{rank} nas demais. A distância entre a curva ingênua e as outras duas é grande ao longo de todo o intervalo testado, o que reforça o efeito do balanceamento das árvores. Entre união por rank e Tarjan, o total de acessos não cai na mesma proporção do tempo de execução: a compressão de caminho acrescenta escritas ao achatar a estrutura, mas compensa isso ao reduzir a profundidade das buscas seguintes.

\begin{figure}[ht]
\centering
\begin{tikzpicture}
\begin{axis}[
  width=0.88\textwidth,
  height=6.8cm,
  xlabel={$n$},
  ylabel={Acessos totais (leituras + escritas)},
  xmode=log,
  log basis x=2,
  ymode=log,
  legend style={at={(0.03,0.97)},anchor=north west,font=\footnotesize},
  grid=major,
]
\addplot+[thick,mark=*,mark size=1.1pt,color=red!75!black] table[x=n,y=naive_acc,col sep=tab] {data/acessos_por_n.dat};
\addplot+[thick,mark=square*,mark size=1.1pt,color=blue!70!black] table[x=n,y=union_rank_acc,col sep=tab] {data/acessos_por_n.dat};
\addplot+[thick,mark=triangle*,mark size=1.3pt,color=green!50!black] table[x=n,y=tarjan_acc,col sep=tab] {data/acessos_por_n.dat};
\legend{Ingênua,União por rank,Rank + compressão}
\end{axis}
\end{tikzpicture}
\caption{Total de acessos instrumentados na mesma sequência de operações (\texttt{parent} em todas as variantes; \texttt{rank} adicional nas duas com rank).}
\label{fig:acessos}
\end{figure}

\subsection{Tabela-resumo}

A Tabela~\ref{tab:resumo} consolida valores nos extremos do conjunto de valores de $n$ utilizados, permitindo comparar ordens de grandeza de tempo e de acessos sem depender da leitura apenas das curvas.

\begin{table}[ht]
\centering
\caption{Resumo numérico para o menor e o maior $n$ do experimento ($m=20n$, 25\% \texttt{Union}). Tempos em segundos; acessos conforme contadores do código.}
\label{tab:resumo}
\begin{tabular}{c c l r r}
\hline
$n$ & $m$ & Variante & Tempo (s) & Acessos \\
\hline
500 & 10000 & Ingênua & 0.000822 & 567\,178 \\
500 & 10000 & União por rank & 0.000110 & 36\,514 \\
500 & 10000 & Rank + compressão & 0.000071 & 48\,507 \\
64000 & 1280000 & Ingênua & 34.431284 & 9\,715\,677\,951 \\
64000 & 1280000 & União por rank & 0.016787 & 4\,804\,277 \\
64000 & 1280000 & Rank + compressão & 0.009553 & 6\,191\,669 \\
\hline

\end{tabular}
\end{table}

\subsection{Síntese}

No cenário aleatório adotado neste trabalho, as três implementações se separam de forma clara. A variante ingênua apresenta crescimento muito mais acentuado, enquanto união por rank e compressão de caminho mantêm curvas bem mais contidas. Esses resultados são compatíveis com a discussão teórica da Seção~\ref{sec:fundamentacao}, embora o experimento, por si só, não substitua a prova formal das taxas assintóticas.

\section{Conclusão}
\label{sec:conclusao}

Os experimentos realizados mostram que as heurísticas de união por rank e compressão de caminho têm efeito direto no desempenho de estruturas de conjuntos disjuntos quando submetidas a sequências mistas de \texttt{Find} e \texttt{Union}. A implementação ingênua apresentou crescimento muito mais acentuado, enquanto as variantes otimizadas mantiveram comportamento mais estável, sobretudo a combinação completa associada a Tarjan, em acordo com a análise amortizada em termos de $\alpha(n)$ \cite{tarjan1975efficiency}.

Esse resultado ajuda a explicar por que algoritmos que fazem muitas consultas de conectividade, como Kruskal em grafos maiores, se beneficiam tanto dessas otimizações. Como continuidade natural deste estudo, podem-se investigar cenários com outras proporções entre uniões e consultas, além de integrar o DSU a uma implementação completa de Kruskal para observar, no programa inteiro, o peso relativo da estrutura de dados em comparação com a etapa de ordenação das arestas.


\bibliographystyle{sbc}
\bibliography{sbc-template}

\appendix
\section*{Apêndice}
\section{Reprodução dos experimentos}
\label{apendice:reproducao}

O código-fonte deste trabalho está disponível em \url{https://github.com/Lucas-Lutti/rabalho-Pr-tico---An-lise-de-Complexidade.git}. A partir da raiz do repositório do código (pasta que contém \texttt{Makefile}, \texttt{src/} e \texttt{artigo/}), os resultados do artigo podem ser reproduzidos com os comandos a seguir:

\begin{itemize}
  \item \textbf{Compilar o benchmark:} \texttt{make}
  \item \textbf{Gerar \texttt{results/benchmark.csv}:} \texttt{./bin/benchmark --all}
  \item \textbf{Atualizar dados para os gráficos e a tabela no \LaTeX{}:} \texttt{python3 scripts/export\_latex\_plot\_data.py}
  \item \textbf{Compilar o PDF do artigo:} \texttt{cd artigo \&\& ./compile\_pdf.sh}
\end{itemize}

\noindent Parâmetros fixos do modo \texttt{--all} (vide \texttt{src/experimentos.cpp}): valores de $n$ $\{500,1000,2000,4000,8000,16000,32000,64000\}$; $m=20n$; fração de operações \texttt{Union} igual a $0{,}25$; semente da instância $i$-ésima dada por $42+7919i$ (com $i=0$ para $n=500$); aquecimento com \texttt{DsuTarjan} sobre a mesma sequência antes das medições.

\section{Formato do arquivo \texttt{results/benchmark.csv}}
\label{apendice:csv}

Cada linha após o cabeçalho corresponde a uma variante e a um par $(n,m)$. Colunas:

\begin{itemize}
  \item \texttt{variant}: \texttt{naive}, \texttt{union\_rank} ou \texttt{tarjan};
  \item \texttt{n}, \texttt{m}: tamanho do universo e número de operações;
  \item \texttt{seed}: semente usada na geração da sequência;
  \item \texttt{time\_ns}: tempo total em nanosegundos para replay da sequência;
  \item \texttt{reads}, \texttt{writes}: contagem instrumentada de leituras e escritas em \texttt{parent}/\texttt{rank};
  \item \texttt{total\_accesses}: soma \texttt{reads + writes}.
\end{itemize}

\section{Instrumentação de acessos}
\label{apendice:stats}

A estrutura \texttt{AccessStats} acumula leituras e escritas registradas pelos métodos \texttt{parent\_at}, \texttt{parent\_set}, \texttt{rank\_at} e \texttt{rank\_set}. A variante ingênua utiliza apenas \texttt{parent}; as outras duas também contabilizam acessos a \texttt{rank}.

\begingroup
\lstset{
  language=C++,
  basicstyle=\ttfamily\footnotesize,
  keywordstyle=\color{jpurple}\bfseries,
  commentstyle=\color{verde},
  numbers=left,
  numberstyle=\tiny,
  breaklines=true,
  tabsize=2,
  xleftmargin=1.5em,
  frame=single
}
\lstinputlisting[caption={\texttt{include/estatisticas\_acesso.hpp} (contadores).},label=lst:access]{../include/estatisticas_acesso.hpp}
\endgroup

\section{Geração reproduzível das operações}
\label{apendice:ops}

\begingroup
\lstset{
  language=C++,
  basicstyle=\ttfamily\footnotesize,
  keywordstyle=\color{jpurple}\bfseries,
  commentstyle=\color{verde},
  numbers=left,
  numberstyle=\tiny,
  breaklines=true,
  tabsize=2,
  xleftmargin=1.5em,
  frame=single
}
\lstinputlisting[caption={Trecho de \texttt{src/operacoes.cpp}: sequência pseudoaleatória com \texttt{std::mt19937}.},label=lst:ops,firstline=7,lastline=28]{../src/operacoes.cpp}
\endgroup

\section{Implementações \texttt{Find} e \texttt{Union} (ingênua e Tarjan)}
\label{apendice:dsu}

A Listagem~\ref{lst:naive} mostra a variante sem rank nem compressão. A Listagem~\ref{lst:tarjan} ilustra compressão de caminho em duas passagens e, em seguida, a rotina de \texttt{unite} por rank utilizada também na variante ``união por rank'' (sem o segundo laço de compressão em \texttt{find}).

\begingroup
\lstset{
  language=C++,
  basicstyle=\ttfamily\footnotesize,
  keywordstyle=\color{jpurple}\bfseries,
  commentstyle=\color{verde},
  numbers=left,
  numberstyle=\tiny,
  breaklines=true,
  tabsize=2,
  xleftmargin=1.5em,
  frame=single
}
\lstinputlisting[caption={\texttt{src/dsu\_ingenuo.cpp}: \texttt{find} e \texttt{unite}.},label=lst:naive,firstline=32,lastline=50]{../src/dsu_ingenuo.cpp}
\lstinputlisting[caption={\texttt{src/dsu\_rank\_com\_compressao.cpp}: \texttt{find} com compressão e \texttt{unite} por rank.},label=lst:tarjan,firstline=48,lastline=82]{../src/dsu_rank_com_compressao.cpp}
\endgroup

\noindent\textit{Nota:} as listagens leem arquivos em \texttt{artigo/include/} e \texttt{artigo/src/}, mantidos junto ao projeto \LaTeX{} para tornar a compilação do artigo autocontida. Após alterar o código-fonte principal na raiz do repositório, essas cópias devem ser atualizadas antes de recompilar o PDF.


\end{document}
